\chapter*{Lời Mở Đầu}
\addcontentsline{toc}{chapter}{Lời Mở Đầu}
\markboth{Lời Mở Đầu}{Lời Mở Đầu}

\begin{truyen}{Lời Mở Đầu}{Opening Words}
\chuhoa{C}{ó ba thứ trong lớp học mà không ai dạy thẳng nhưng ai cũng học được nếu chú ý: nghệ thuật mắng mà không làm người kia tức, nghệ thuật cà khịa mà người bị khịa phải tự nhìn lại mình, và nghệ thuật hỏi ngây thơ đến mức câu trả lời của thầy cô trở thành bài học đời thật.}
\textit{(There are three things no one teaches directly in a classroom yet everyone can learn if they pay attention: the art of scolding without provoking anger, the art of roasting someone so they reflect on themselves, and the art of asking innocent questions that turn a teacher's reply into a real-life lesson.)}

Quyển này ghi lại đúng ba thứ đó. Không có chuyện cổ tích, không có nhân vật hoàn hảo. Chỉ có lớp học thật, giọng thầy thật, câu hỏi thật của học sinh thật --- và những câu đáp mà nếu nghe xong không nghĩ ngợi thêm một chút thì đọc lại lần nữa.
\textit{(This volume records exactly those three things. No fairy tales, no perfect characters. Just real classrooms, real teacher voices, real student questions — and replies that, if they don't make you pause to think, deserve a second read.)}

Ba chủ đề, mỗi chủ đề hướng đến 1000 câu, chia thành nhiều quyển. Đây là Phần I. Những câu tiếp theo sẽ đến, miễn là các em vẫn còn hỏi những câu như thế này.
\textit{(Three themes, each aiming for 1000 entries, spread across many volumes. This is Part I. The next entries will come — as long as students keep asking questions like these.)}
\end{truyen}

\begin{baihoc}
Câu nói đùa mà thoáng nghe thấy đau — đó là câu nói đang chạm đúng chỗ cần chạm.\\
\textit{(A joke that stings slightly when you hear it — that is the one hitting exactly where it should.)}
\end{baihoc}

\ngancach

\begin{tuhocnhanh}
\textbf{Cách đọc sách này:}\\[4pt]
\begin{itemize}[leftmargin=*]
  \item Đọc câu tiếng Việt trước. Nghĩ xem mình có nhận ra ai trong đó không.
  \item Đọc câu tiếng Anh. Xem sắc thái có thay đổi không.
  \item Nếu một câu làm mình hơi khó chịu, đừng bỏ qua --- đọc lại và hỏi tại sao.
\end{itemize}
\textit{(How to read this book: Read the Vietnamese line first. See if you recognize anyone in it. Read the English. Notice whether the tone shifts. If a line makes you uncomfortable, don't skip it --- reread and ask yourself why.)}
\end{tuhocnhanh}
